\section{Conclusion}

Ce laboratoire nous a permis de programmer un contrôleur CAN MCP2515 et
d'observer son comportement réel à l'oscilloscope.

Plusieurs points sont à retenir. D'abord, le protocole CAN est robuste :
si personne n'acquitte le message, l'émetteur le retransmet sans cesse
jusqu'à obtenir une réponse. Ce comportement est utile pour la fiabilité,
mais peut bloquer le bus en cas de panne.

Ensuite, la latence dépend fortement de la charge. Sur un bus libre,
chaque message passe rapidement et la gigue reste faible. Dès qu'un autre
flux est présent, certains messages doivent attendre, et la
gigue augmente.

Enfin, l'essai avec un signal sinusoïdal illustre clairement la
saturation : à \num{200}~kbps, tant qu'un seul flux circule la sortie
suit correctement l'entrée, mais en ajoutant les flux un par un la qualité
se dégrade jusqu'à devenir totalement inexploitable avec quatre flux à
\num{50}~Hz. Le bus est alors saturé : les messages arrivent trop tard
pour reconstruire le signal.

Le choix du débit n'est donc pas le seul critère lors du dimensionnement
d'une liaison CAN. Le nombre de nœuds, la fréquence d'émission et le
type de signal transmis jouent un rôle au moins aussi important pour
garantir une communication exploitable en pratique.
